% Part: normal-modal-logic
% Chapter: sequent-calculus
% Section: rules-for-K

\documentclass[../../../include/open-logic-section]{subfiles}

\begin{document}

\olfileid{nml}{seq}{rul}

\olsection{\Log{K} માટેના નિયમો}

સામાન્ય વિધાનસંબંધી સંયોજકો માટેના નિયમો સામાન્ય સિક્વન્ટ
કલન~$\Log{LK}$ના નિયમો જેવા જ છે. સ્વયંસિદ્ધિઓ પણ એ જ છે:
$!A \Sequent !A$ સ્વરૂપનું કોઈપણ સિક્વન્ટ સ્વયંસિદ્ધિ ગણાય છે.

મોડલ સંક્રિયક\iftag{prvBox}{\iftag{prvDiamond}{ો~$\Box$
અને}{~$\Box$}}{}\iftag{prvDiamond}{~$\Diamond$}{} માટે નીચેના
વધારાના \iftag{notprvBox,notprvDiamond}{નિયમ}{નિયમો} છે:
\[
  \iftag{prvBox}{\iftag{prvDiamond}{% <> and [] primitive
    \Axiom$\Gamma \fCenter \Delta, !A$
    \RightLabel{$\Box$}
    \UnaryInf$\Box\Gamma \fCenter \Diamond\Delta, \Box!A$
    \DisplayProof  
    \qquad
    \Axiom$!A, \Gamma \fCenter \Delta$
    \RightLabel{$\Diamond$}
    \UnaryInf$\Diamond!A, \Box\Gamma \fCenter \Diamond\Delta$
    \DisplayProof
}{% only Box primitive
\Axiom$\Gamma \fCenter !A$
\RightLabel{$\Box$}
\UnaryInf$\Box\Gamma \fCenter \Box!A$
\DisplayProof
}}{% only <> primitive
  \Axiom$!A \fCenter \Delta$
  \RightLabel{$\Diamond$}
  \UnaryInf$\Diamond!A \fCenter \Diamond\Delta$
  \DisplayProof
}
\]
અહીં \iftag{prvBox}{$\Box\Gamma$ એટલે~$\Gamma$ના દરેક
!!{formula} પહેલાં $\Box$ મૂકવાથી~$\Gamma$માંથી મળતી
!!{formula}sની
શ્રેણી\iftag{prvDiamond}{ અને
}{}}{}\iftag{prvDiamond}{$\Diamond\Delta$ એટલે~$\Delta$ના દરેક
!!{formula} પહેલાં $\Diamond$ મૂકવાથી~$\Delta$માંથી મળતી
!!{formula}sની
શ્રેણી}{}. \iftag{notprvDiamond}{$\Box$ નિયમના પૂર્વપક્ષની
જમણી બાજુએ વધુમાં વધુ એક જ !!{formula}~$!A$ હોવું જોઈએ. }{}%
\iftag{notprvBox}{$\Diamond$ નિયમના પૂર્વપક્ષની ડાબી
બાજુએ વધુમાં વધુ એક જ !!{formula}~$!A$ હોવું જોઈએ. }{}%
\iftag{prvBox}{$\Gamma$\iftag{prvDiamond}{~અને~}{}}{}\iftag{prvDiamond}{$\Delta$}{}
ખાલી હોઈ શકે; ત્યારે નિષ્કર્ષ સિક્વન્ટનો અનુરૂપ ભાગ
\iftag{prvBox}{$\Box\Gamma$\iftag{prvDiamond}{~અને~}{}}{}\iftag{prvDiamond}{$\Diamond\Delta$}{}
પણ ખાલી હોય છે.

\iftag{prvBox}{જમણી બાજુ $\Box$\iftag{prvDiamond}{ અને
}{}}{}\iftag{prvDiamond}{ડાબી બાજુ $\Diamond$}{} માત્ર એક જ
!!{formula}~$!A$ પર મૂકવાની મર્યાદા જરૂરી છે. જો
\iftag{prvBox}{જમણી બાજુ ગમે તેટલાં !!{formula}s પર $\Box$
મૂકવાની\iftag{prvDiamond}{ અથવા }{}}{}\iftag{prvDiamond}{ડાબી
બાજુ ગમે તેટલાં !!{formula}s પર $\Diamond$ મૂકવાની}{}
છૂટ આપીએ, તો નીચેનું !!{derive} કરી શકીએ:
  \[
    \iftag{prvBox}{
      \Axiom$!A \fCenter !A$
      \RightLabel{\RightR{\lnot}}
      \UnaryInf$\fCenter !A, \lnot !A$
      \RightLabel{$\Box*$}
      \UnaryInf$\fCenter \Box!A, \Box\lnot !A$
      \doubleLine
      \RightLabel{\RightR{\lor}}
      \UnaryInf$\fCenter \Box!A \lor \Box \lnot !A$
      \DisplayProof}{}
    \iftag{notprvBox,notprvDiamond}{}{\qquad}
    \iftag{prvDiamond}{
      \Axiom$!A \fCenter !A$
      \RightLabel{\LeftR{\lnot}}
      \UnaryInf$\lnot !A, !A \fCenter$
      \RightLabel{$\Diamond*$}
      \UnaryInf$\Diamond\lnot !A,\Diamond!A \fCenter $
      \RightLabel{\RightR{\lnot}}
      \UnaryInf$\Diamond!A \fCenter \lnot\Diamond\lnot !A$
      \RightLabel{\RightR{\lif}}
      \UnaryInf$ \fCenter \Diamond!A \lif \lnot\Diamond\lnot !A$
      \DisplayProof}{}
  \]
પરંતુ \iftag{prvBox}{$\Box!A \lor \Box \lnot !A$\iftag{prvDiamond}{ અને
$\Diamond!A \lif \lnot\Diamond\lnot !A$ બંને}{}}{$\Diamond!A \lif
\lnot\Diamond\lnot !A$}~$\Log{K}$માં માન્ય નથી.

જો પૂર્વપક્ષમાં~$!A$ ઉપરાંત બાજુનાં સૂત્રો રાખીએ અને
\iftag{prvBox}{$\Box$ નિયમને જમણી બાજુ માત્ર~$!A$ પર
$\Box$ મૂકવા દઈએ\iftag{prvDiamond}{ અથવા }{}}{}\iftag{prvDiamond}{$\Diamond$
નિયમને ડાબી બાજુ માત્ર~$!A$ પર $\Diamond$ મૂકવા દઈએ}{}
(પણ બાજુનાં સૂત્રોને યથાવત્ રાખીએ), તો નીચેનું
!!{derive} કરી શકીએ:
\[
  \iftag{prvBox}{
    \Axiom$!A \fCenter !A$
    \RightLabel{\RightR{\lnot}}
    \UnaryInf$\fCenter !A, \lnot !A$
    \RightLabel{\RightR{\Exchange}}
    \UnaryInf$\fCenter \lnot !A, !A$
    \RightLabel{$\Box*$}
    \UnaryInf$\fCenter \lnot!A, \Box !A$
    \doubleLine
    \RightLabel{\RightR{\lor}}
    \UnaryInf$\fCenter \lnot!A \lor \Box !A$
    \DisplayProof}{}
  \iftag{notprvBox,notprvDiamond}{}{\qquad}
  \iftag{prvDiamond}{
    \Axiom$!A \fCenter !A$
    \RightLabel{\LeftR{\lnot}}
    \UnaryInf$\lnot !A, !A \fCenter$
    \RightLabel{$\Diamond*$}
    \UnaryInf$\Diamond\lnot !A, !A \fCenter $
    \RightLabel{\RightR{\lnot}}
    \UnaryInf$!A \fCenter \lnot\Diamond\lnot !A$
    \RightLabel{\RightR{\lif}}
    \UnaryInf$ \fCenter !A \lif \lnot\Diamond\lnot !A$
    \DisplayProof}{}
\]
પરંતુ \iftag{prvBox}{$\lnot!A \lor \Box !A$ (જે $!A
\lif \Box!A$ને સમકક્ષ છે)\iftag{prvDiamond}{ અને
$!A \lif \lnot\Diamond\lnot !A$ બંને}{}}{$!A \lif
\lnot\Diamond\lnot !A$}~$\Log{K}$માં માન્ય નથી.

\end{document}
